\ifx\allfiles\undefined
\documentclass[12pt, a4paper, oneside, UTF8]{ctexbook}
\def\configPath{../config}
\input{\configPath/config}
\begin{document}
\else
\fi

\chapter{实分析}
\section*{第一小部分：一元实变量函数的Lebesgue积分}
\section*{集合、映射与关系的预备知识}

\section{实数集：集合、序列与函数}

\section{Lebesgue测度}


\section{Lebesgue可测函数}
\begin{proposition}\label{proposhifenxi1}
	令函数$f$具有可测定义域$E$,则下面叙述等价
	\begin{change}
		\item $\forall c \in R,\{x\in E| f(x)>c\} \in \mathscr{U}$ ,$\mathscr{U}$为可测集族  
		\item $\forall c \in R,\{x\in E| f(x)\geqslant c\} \in \mathscr{U}$
		\item $\forall c \in R,\{x\in E| f(x)<c\} \in \mathscr{U}$
		\item $\forall c \in R,\{x\in E| f(x)\leqslant  c\} \in \mathscr{U}$
	\end{change}
\end{proposition}

\begin{defn}
	定义在$E$上的扩张的实值函数$f$称为是Lebesgue可测的，若它的定义域$E$是可测的且满足\cref{proposhifenxi1}的四个命题之一的。
\end{defn}

\begin{proposition}
	令$f$为可测集合$E$上的实值函数，则函数$f$是可测的当且仅当对每个开集 %$\mathscr{O}$，
	$\mathcal{O}$,$\mathcal{O}$在$f$下的原象为$f^{-1}(\mathcal{O})=\{x \in E|f(x) \in \mathcal{O}\}$是可测的。
\end{proposition}		

\begin{proposition}
	\begin{change}
		\item 在可测集定义域上连续的实值函数是可测的。
		\item 定义在区间上的可测函数是可测的。
		\item 令$f$和$g$为$E$上a.e.有限的可测函数，则线性（$\alpha f + \beta g$）和（$f\cdot g$）也可测。
		\item  ui
	\end{change}
\end{proposition}

\begin{proposition}
	\begin{enumerate}[label=\Roman*]
		 令$f$为$	E$上扩充的实值函数。
		\item  若$f$为可测的，而在$E$上$f=g \ a.e.$,则$g$在$E$上可测。
		\item 对于$E$的可测子集$D$,$f$是可测的当且仅当在$D$和$E~D$上的限制是可测的。
	\end{enumerate}
\end{proposition}



\section{Lebesgue积分}



\section{微分与积分}

\section{$L^P$空间：完备性与逼近}



\section{$L^P$空间：对偶与弱收敛}

\ifx\allfiles\undefined
\end{document}
\fi